\documentclass[journal]{IEEEtran}

\usepackage{cite}
\usepackage{amsmath,amssymb,amsfonts}
\usepackage{algorithmic}
\usepackage{graphicx}
\usepackage{textcomp}
\usepackage{xcolor}
\usepackage{booktabs}
\usepackage{multirow}
\usepackage{hyperref}

% ---------------------------------------------------------------------------
% Draft-only helper: visually flags unfinished content so nothing pending is
% mistaken for a finished claim. Remove this package/command before final
% submission (do NOT leave \todo markers in a submitted manuscript).
% ---------------------------------------------------------------------------
\newcommand{\todo}[1]{\textcolor{red}{\textbf{[TODO: #1]}}}

\begin{document}

\title{DH$^2$A-KT: Dynamic Heterogeneous Hypergraph and Agentic Knowledge
Tracing with Leakage-Controlled Graph Construction}

\author{Tuan~Dao~Minh, Khanh-Trinh~Nguyen, Duong~Nguyen~Tien, Quoc~Khanh~Ngo,
Van-Hau~Nguyen, and~Le~Hoang~Son
\thanks{T. Dao Minh, K.-T. Nguyen, D. N. Tien, and V.-H. Nguyen are with the
Department of Quality Assurance and Testing, Hung Yen University of
Technology and Education, Hung Yen, Vietnam (e-mail: tuanymc@utehy.edu.vn;
trinhnk@utehy.edu.vn; duongnt@utehy.edu.vn; nvhau66@gmail.com).}
\thanks{Q. K. Ngo and L. H. Son are with the Artificial Intelligence Research
Center, VNU Information Technology Institute, Vietnam National University,
Hanoi, Vietnam (e-mail: quockhanhngo.official@gmail.com; sonlh@vnu.edu.vn).
L. H. Son is the corresponding author.
\todo{Manuscript received/revised dates -- fill at submission time.}}}

\markboth{IEEE Transactions on Learning Technologies}%
{Dao \MakeLowercase{\textit{et al.}}: DH$^2$A-KT}

\maketitle

\begin{abstract}
Graph-augmented knowledge tracing (KT) models student mastery by propagating
information over auxiliary concept graphs, but two gaps limit current
practice. First, most graph-KT architectures represent only pairwise
knowledge-component relations and omit heterogeneous entities -- teachers,
hints, forum discussions, instructional videos -- that plausibly shape a
learner's trajectory. Second, when graph-augmented KT models are evaluated
without auditing whether the auxiliary graph was itself built without
leakage, reported accuracy gains are difficult to attribute to genuine
structural signal versus held-out information reaching the model through the
graph channel. We present DH$^2$A-KT, a two-tier architecture that addresses
both gaps while remaining computationally tractable on a single consumer
GPU. Tier 1 (DH$^2$-KT) extends a leakage-controlled, train-only concept
graph construction and audit protocol -- released with a companion paper,
P0 -- from pairwise prerequisite/similarity edges to heterogeneous
hyperedges (session, concept-prerequisite, discussion-thread,
teacher-intervention), and adds a propensity-score causal-effect layer to
separate genuine intervention effects from spurious correlation. We extend
P0's four-diagnostic leakage taxonomy with three hyperedge-specific leakage
classes -- group-membership, cross-modal, and intervention-timing leakage --
that a pairwise audit cannot detect. Tier 2 (AgentDG-KT) is a
deliberately-scoped pilot in which a small team of LLM agents
(Diagnostician, Tutor/Hint, Critic) reads Tier 1's frozen output to produce
natural-language explanations, gated by a mandatory Critic agent rather than
benchmarked at the same quantitative scale as Tier 1. We reuse P0's already
audited, leakage-controlled baseline results on XES3G5M, ASSISTments 2012,
and Junyi Academy as the Tier 1 comparison column, avoiding redundant
retraining of nine baseline knowledge-tracing models.
\todo{Replace this sentence with the actual Tier 1/Tier 2 headline result
once M5/M6/M9 (see docs/execution-plan.md) are complete -- do not
pre-announce a number before it exists.}
\end{abstract}

\begin{IEEEkeywords}
Knowledge tracing, hypergraph neural networks, heterogeneous graphs, large
language model agents, causal inference, data leakage, educational data
mining.
\end{IEEEkeywords}

% =============================================================================
\section{Introduction}
% =============================================================================

Knowledge tracing (KT) estimates a learner's latent mastery of knowledge
components (KCs) from interaction history, supporting adaptive practice,
intelligent tutoring, and curriculum recommendation. The field has
progressed from probabilistic and sequence models
\cite{piech2015dkt} to graph-augmented models that propagate information
over auxiliary concept graphs \cite{nakagawa2019gkt,yang2021gikt}, and, most
recently, toward hypergraph representations that capture higher-order
relations among knowledge components \cite{thmn2025,hgkt2025,htkt2025} and
dynamic, temporally-evolving graph structure \cite{dygkt2024}.

This paper addresses these three gaps directly, building on -- rather than
replacing -- \cite{daominh2026p0}'s leakage-control protocol. Tier~1
(Sections~\ref{sec:tier1-graph}--\ref{sec:causal}) extends that protocol's
train-only, audited concept-prerequisite edges into heterogeneous
hyperedges and adds the causal-effect layer the first and third gaps call
for; Section~\ref{sec:leakage-ext} extends its four-diagnostic leakage
audit to the hyperedge setting the second gap identifies as unaudited.
Tier~2 (Section~\ref{sec:tier2}) is a deliberately-scoped pilot showing that
an LLM-agent interpretability layer can read Tier~1's frozen, audited
output without reopening the leakage question Tier~1 already closed.
\todo{Once M5/M6/M9 (docs/execution-plan.md) produce results, replace this
paragraph's forward-looking framing with a results-grounded framing and add
the C6 headline-finding sentence below -- do not claim a gain exists before
it is measured.}

\subsection{Motivation}
Three concrete gaps motivate this work:

\begin{itemize}
  \item \textbf{Heterogeneous relations are absent from graph-KT.}
  Hypergraph-KT models \cite{thmn2025,hgkt2025,htkt2025} group multiple
  knowledge components into a single hyperedge, but the hyperedge always
  contains only KC nodes. Teacher intervention, hint content, peer
  discussion, and instructional video -- entities with documented influence
  on learning outcomes -- are not represented as graph nodes in any
  published graph-KT model we are aware of.
  \item \textbf{Graph provenance is rarely audited.} \cite{daominh2026p0}
  shows that pooling a knowledge-component graph from the full interaction
  log before a learner-based train/test split lets held-out information
  reach a graph-augmented KT model even when the sequence optimizer itself
  trains only on the training fold. \cite{daominh2026p0}'s protocol audits
  this channel for \emph{pairwise} prerequisite/similarity edges; no
  published audit protocol currently covers \emph{hyperedges} spanning
  heterogeneous entity types.
  \item \textbf{Causal claims about pedagogical intervention are rare.} Most
  graph-KT models report predictive accuracy only, without attempting to
  separate a hint or teacher action's causal effect on subsequent
  correctness from confounded correlation.
\end{itemize}

\subsection{Contributions}
\begin{itemize}
  \item \textbf{C1.} A heterogeneous hypergraph construction procedure
  (Section~\ref{sec:tier1-graph}) that groups \cite{daominh2026p0}'s
  already leakage-audited, acyclic prerequisite edges into
  concept-prerequisite hyperedges, and defines (interface-complete;
  data-blocked, see Section~\ref{sec:limitations}) session,
  discussion-thread, and teacher-intervention hyperedge types.
  \item \textbf{C2.} A hyperedge-level extension of \cite{daominh2026p0}'s
  four-diagnostic leakage taxonomy, adding three classes --
  group-membership, cross-modal, and intervention-timing leakage -- that a
  pairwise audit structurally cannot detect (Section~\ref{sec:leakage-ext}).
  \item \textbf{C3.} A causal-effect layer (Section~\ref{sec:causal})
  estimating the treatment effect of a hint or teacher intervention via
  inverse-propensity weighting, as a design-target starting point before a
  more elaborate causal architecture is attempted.
  \item \textbf{C4.} Direct reuse of \cite{daominh2026p0}'s published,
  leakage-controlled baseline results (BKT, DKT, AKT, simpleKT, GKT, GIKT,
  SKT, DyGKT, DGEKT on XES3G5M/ASSISTments~2012/Junyi Academy) as the Tier~1
  comparison column, avoiding redundant baseline retraining under matched
  protocol and split seeds.
  \item \textbf{C5.} A deliberately-scoped Tier~2 pilot (AgentDG-KT,
  Section~\ref{sec:tier2}) demonstrating that an LLM-agent interpretability
  layer reading Tier~1's frozen output -- rather than re-inferring knowledge
  state from scratch -- is tractable on a single consumer GPU, gated by a
  mandatory Critic agent.
\end{itemize}

\todo{Once results exist, add a C6 sentence stating the headline empirical
finding (or its absence/negative result) -- P0's own Introduction models
this pattern (see its C5) and it should be mirrored here honestly, including
if Tier 1's hypergraph gain turns out to be small or backbone-specific.}

% =============================================================================
\section{Related Work}
% =============================================================================

\subsection{Graph- and hypergraph-based knowledge tracing}
Early graph-augmented KT models propagate information over a single,
pairwise concept graph: \cite{nakagawa2019gkt} layer a graph neural network
over a knowledge-component graph on top of a sequence model, and
\cite{yang2021gikt} let exercises and concepts recursively re-rank one
another through a bipartite exercise-concept graph. \cite{dygkt2024} extends
this line to a temporally-evolving graph that updates edges as new
interactions arrive rather than freezing the graph after construction. A
second line moves from pairwise edges to hyperedges spanning more than two
knowledge components at once: \cite{thmn2025} proposes a temporal hypergraph
memory network, \cite{hgkt2025} adds a dual-gated mechanism controlling how
much a hyperedge's aggregated signal influences a given prediction, and
\cite{htkt2025} replaces graph convolution with a hypergraph transformer. A
third line adds heterogeneity in what a node \emph{represents} rather than
how many nodes an edge spans: \cite{psykt2024} augments the concept graph
with psychological-factor nodes, and \cite{hhskt2023} builds a heterogeneous
graph directly over learner-question interactions. Across all three lines,
every published hyperedge or heterogeneous node we are aware of is still
drawn from a closed set of KC-, exercise-, or learner-adjacent entities; none
represent a Teacher intervention, Hint content, Forum discussion, or
instructional Video as a first-class member inside the same hyperedge as the
knowledge components it plausibly affects -- the gap C1
(Section~\ref{sec:tier1-graph}) addresses.

\subsection{LLM agents in education and graph-augmented LLM agents}
\cite{llmedu2025} surveys the growing use of LLM agents across tutoring,
feedback generation, and curriculum planning in education, and
\cite{agentcat2026} instantiates a multi-agent LLM pipeline that simulates
computerized adaptive testing end to end, with agents standing in for both
the test administrator and the examinee. Outside education,
\cite{graphaugmentedllm2025} surveys how LLM agents can read from and write
to an external graph as working memory, and \cite{graphrag2025} applies a
graph-retrieval-augmented-generation pattern to recommend personalized
learning paths over a dual knowledge-structure graph. AgentDG-KT
(Section~\ref{sec:tier2}) differs from this pattern in one specific way:
none of the above agents are gated by a design in which the graph -- and the
quantitative model that reads it -- is frozen before the agent acts, with
the agent limited to producing a natural-language artifact (rationale,
hint, session hyperedge) that a separate Critic agent must approve before it
is treated as anything more than a proposal. We adopt this
frozen-core-plus-pilot-agent split specifically because a full multi-agent
LLM pipeline that re-infers knowledge state end to end over millions of
interactions is not tractable on a single consumer GPU
(Section~\ref{sec:tier2}), not because we believe agentic re-inference is
undesirable in principle.

\subsection{Leakage-controlled evaluation in knowledge tracing}
\cite{daominh2026p0} is the companion paper this work builds on and reuses
directly (Section~\ref{sec:reuse}); we summarize its protocol here so the
rest of this paper is self-contained. \cite{daominh2026p0} observes that a
knowledge-component graph is often constructed by pooling co-occurrence or
similarity statistics over the \emph{entire} interaction log before a
learner-based train/test split is applied, which lets held-out information
reach a graph-augmented KT model through the graph channel even when the
sequence model itself trains only on the training fold. Its audit protocol
reports four scalar diagnostics computed on a train-only-constructed graph
-- an edge-composition-ratio flag ($\mathrm{ECR}^{\mathrm{flag}}$), an
edge-composition overlap ratio ($\mathrm{ECR}^{\mathrm{overlap}}$), an
edge-outcome correlation ($|\rho|$), and a train/test boundary-mismatch rate
($\mathrm{TBMR}$) -- alongside a DAG audit with cycle pruning, a cold-start
stratification of knowledge components into four frequency strata, a DAG
Disruption Rate (DDR) manipulation check under five structured augmentation
operators, and a ground-truth cross-validation against an expert-curated
prerequisite DAG on Junyi Academy. Across XES3G5M, ASSISTments~2012, and
Junyi Academy, \cite{daominh2026p0} reports that train-only graph
construction shifts baseline AUC by at most 0.003 relative to full-log
construction -- evidence that leakage control, in their setting, does not by
itself manufacture the reported gains, and a result we rely on directly when
interpreting any Tier~1 gain in this paper (Section~\ref{sec:reuse}).
Section~\ref{sec:leakage-ext} extends this four-diagnostic protocol from
pairwise edges to heterogeneous hyperedges.

\subsection{Causal inference in educational data mining}
Causal analysis in education has mostly targeted two separate questions:
estimating the effect of a known treatment, and discovering an unknown
causal structure among skills. \cite{kaur2019causal} addresses the first:
given historical course-enrollment and grade data, they estimate the
average treatment effect of taking one course before another using matching
and regression -- directly analogous in spirit to the propensity-score
estimator we use for hint/teacher-intervention effects
(Section~\ref{sec:causal}), though on a coarser course-level treatment
rather than a step-level pedagogical action. \cite{kumar2023causal}
addresses the second: they learn a latent causal ordering and dependency
structure among knowledge-component skills end to end inside a modified DKT
model, evaluated on the NeurIPS 2022 Challenge on Causal Insights for
Learning Paths in Education. \cite{minn2022interpretable} sits between the
two, using a causal-relations-constrained student model to make KT
predictions more interpretable without a separate treatment-effect
estimation step. Our causal-effect layer (Section~\ref{sec:causal}) is
closest in spirit to \cite{kaur2019causal} -- a treatment-effect estimator
rather than a causal-discovery method -- but applies it at the level of
individual hint/teacher-intervention hyperedges rather than curriculum-level
course sequencing, and explicitly reports the propensity-overlap diagnostic
that \cite{kaur2019causal} does not discuss, following the
identification-assumption caveat in Section~\ref{sec:limitations}.

% =============================================================================
\section{Preliminaries}
% =============================================================================

Let $\mathcal{D} = \{(u, t, q, c, y)\}$ be a KT interaction log with learner
$u$, timestamp $t$, item $q$, knowledge component $c$, and correctness label
$y \in \{0, 1\}$, following \cite{daominh2026p0}'s notation. A learner-based
temporal split $\mathcal{F} = (\mathcal{F}_{\mathrm{train}},
\mathcal{F}_{\mathrm{val}}, \mathcal{F}_{\mathrm{test}})$ assigns each
learner to exactly one fold (ratios $0.7/0.1/0.2$, three folds with split
seeds $42$/$43$/$44$, matching \cite{daominh2026p0}'s protocol so results
remain directly comparable, see Section~\ref{sec:reuse}).

We extend this notation with a \emph{hyperedge} $h = (\kappa,
\{(\tau_i, e_i)\}_{i=1}^{k}, f)$: a kind label $\kappa \in
\{\text{concept-prerequisite}, \text{session}, \text{discussion-thread},
\text{teacher-intervention}\}$, a set of $k$ typed member entities
$(\tau_i, e_i)$ with entity type $\tau_i \in \{\text{student, exercise,
concept, teacher, hint, forum, discussion, video}\}$, and a fold identifier
$f$. A heterogeneous hypergraph $\mathcal{H}_f = (\mathcal{V}_f,
\mathcal{E}_f^{\mathrm{hyper}})$ for fold $f$ is a set of such hyperedges
over the typed node set $\mathcal{V}_f$.

\subsection{Design goal}
For each fold $f$, the goal is a hypergraph $\mathcal{H}_f$ constructed
using only $\mathcal{F}_f^{\mathrm{train}}$ evidence, audited for the
leakage classes in Table~\ref{tab:leakage-taxonomy}
(Section~\ref{sec:leakage-ext}), that a downstream Tier~1 encoder consumes
to predict $y$ for held-out interactions.

% =============================================================================
\section{Tier 1: DH$^2$-KT}
% =============================================================================

\subsection{Heterogeneous hypergraph construction}
\label{sec:tier1-graph}

\todo{Describe build\_concept\_prerequisite\_hyperedges() precisely
(Algorithm~\ref*{alg:concept-hyperedge} below is a first draft matching
dh2a\_kt/hyperedge/construction.py -- keep this in sync with the actual code
as it is implemented per docs/execution-plan.md M1-M2).}

\begin{algorithmic}[1]
\STATE \textbf{Input:} audited, acyclic prerequisite edges
$E_{\mathrm{pre}}^f$ from \cite{daominh2026p0}'s Algorithm~2; chain length
bounds $[\ell_{\min}, \ell_{\max}]$
\STATE \textbf{Output:} concept-prerequisite hyperedge set
$\mathcal{E}^{\mathrm{cprereq}}_f$
\STATE Build directed graph $G$ from $E_{\mathrm{pre}}^f$
\FOR{each root $r$ with in-degree 0 in $G$}
  \FOR{each simple path $p$ from $r$ with $\ell_{\min} \le |p| \le
  \ell_{\max}$}
    \STATE emit hyperedge with members $\{(\text{concept}, c) : c \in p\}$
  \ENDFOR
\ENDFOR
\end{algorithmic}
\label{alg:concept-hyperedge}

Session, discussion-thread, and teacher-intervention hyperedges are
specified with the same interface but are \emph{data-blocked} on all three
core benchmarks: no public KT benchmark studied by \cite{daominh2026p0}
(XES3G5M, ASSISTments~2012, Junyi Academy) ships Hint, Forum, or
Teacher-action logs synchronized to the interaction stream. We resolve this
per hyperedge kind rather than uniformly. Discussion-thread and
teacher-intervention hyperedges have no synchronized public source at all
(Forum and Teacher-action logs are absent from every benchmark we are aware
of); we scope these two hyperedge kinds to the Tier~2 case study
(Section~\ref{sec:tier2}) rather than the Tier~1 quantitative core. Session
hyperedges (Student-Exercise-Concept-Hint(-Video)) have a candidate public
source, ES-KT-24 \cite{esk2024}, which ships educational-game-playing video
alongside interaction logs. \todo{State whether ES-KT-24 wiring
(docs/execution-plan.md M2) succeeded once attempted -- if it does not,
Tier~1's quantitative core remains scoped to concept-prerequisite
hyperedges only, which is still sufficient to train DH$^2$-KT end to end
per docs/idea-D-plan.md section~5.}

\subsection{Hyperedge-level leakage audit}
\label{sec:leakage-ext}

\cite{daominh2026p0} audits five leakage classes at the level of pairwise
edge provenance (structural, edge, target, temporal, cold-start
neighbourhood) using four scalar diagnostics: $\mathrm{ECR}^{\mathrm{flag}}$,
$\mathrm{ECR}^{\mathrm{overlap}}$, $|\rho|$, and $\mathrm{TBMR}$. These
diagnostics are reused unmodified (Section~\ref{sec:reuse}) on the pairwise
projection of a hyperedge set. Table~\ref{tab:leakage-taxonomy} adds three
classes that a pairwise audit cannot express because they require a
\emph{group} of heterogeneous members, not a single edge.

\begin{table}[!t]
\caption{Hyperedge-level leakage taxonomy (extends \cite{daominh2026p0}
Table~1)}
\label{tab:leakage-taxonomy}
\centering
\begin{tabular}{@{}p{2.1cm}p{2.9cm}p{2.4cm}@{}}
\toprule
Class & Source & Diagnostic \\
\midrule
Group-membership & A hyperedge with $\ge 1$ member drawn from a held-out
interaction while other members are train-only & $\mathrm{compute\_group\_}$ $\mathrm{membership\_}$ $\mathrm{leakage}$ \\
Cross-modal & A Forum/Video embedding pooled over content timestamped after
the prediction cutoff & $\mathrm{compute\_cross\_}$ $\mathrm{modal\_leakage}$ \\
Intervention-timing & A Teacher/Hint hyperedge whose content was effectively
chosen using knowledge of the student's later correctness & $\mathrm{compute\_intervention\_}$ $\mathrm{timing\_leakage}$ \\
\bottomrule
\end{tabular}
\end{table}

\todo{Report actual leakage rates here once docs/execution-plan.md M1 runs
on real XES3G5M fold data -- Table~\ref{tab:leakage-taxonomy} currently
documents the diagnostic definitions only, not measured values.}

\subsection{Encoding and hypergraph attention}
\todo{Section pending docs/execution-plan.md M3 (dh2a\_kt/models/dh2\_kt.py
DH2KT.forward() implementation). Do not describe the attention mechanism in
past tense until it exists and has passed the sanity overfit test in M3's
Definition of Done.}

\subsection{Causal-effect layer}
\label{sec:causal}
We estimate the average treatment effect (ATE) of a binary intervention
(e.g., hint given) on next-step correctness via inverse-propensity
weighting, following the identification assumption that intervention
assignment is unconfounded given observed features -- a strong assumption on
observational KT data that we report as a limitation rather than a proven
property (Section~\ref{sec:limitations}). Given confounders $X$, treatment
$T \in \{0,1\}$, and outcome $Y$, the propensity score $\hat{e}(X) =
P(T=1\mid X)$ is estimated by logistic regression, and
\begin{equation}
\widehat{\mathrm{ATE}} = \frac{\sum_i T_i Y_i / \hat{e}(X_i)}{\sum_i T_i /
\hat{e}(X_i)} - \frac{\sum_i (1-T_i) Y_i / (1-\hat{e}(X_i))}{\sum_i (1-T_i) /
(1-\hat{e}(X_i))}.
\label{eq:ate}
\end{equation}
\todo{Report propensity-overlap diagnostics alongside any ATE number, per
the warning already implemented in dh2a\_kt/models/causal\_layer.py -- an
ATE without an overlap check is not publishable on its own.}

% =============================================================================
\section{Tier 2: AgentDG-KT pilot}
\label{sec:tier2}
% =============================================================================

Tier~2 is deliberately scoped as a pilot, not a benchmark at Tier~1's
quantitative rigor (docs/idea-D-plan.md, following the feasibility analysis
in the accompanying project discussion: a full multi-agent LLM re-inference
pipeline over millions of interactions is not tractable on a single
consumer GPU within a realistic project timeline). Three agents operate on
Tier~1's \emph{frozen} output rather than re-inferring knowledge state:

\begin{itemize}
  \item \textbf{Diagnostician Agent} produces a natural-language rationale
  for a given Tier-1 prediction, explicitly instructed not to override the
  given probability.
  \item \textbf{Tutor/Hint Agent} proposes the next hint/exercise and writes
  a new session hyperedge back into the graph (the agentic closed loop).
  \item \textbf{Critic Agent} is a \emph{mandatory} quality gate that flags
  any Diagnostician output contradicting or ignoring the Tier-1 number it
  was meant to explain, before anything is treated as ground truth in a
  case study.
\end{itemize}

\todo{Report pilot results (sample size, hallucination/flag rate, human
evaluation if collected) here once docs/execution-plan.md M8-M9 complete.
Do not report Tier 2 metrics using Tier 1's AUC/ACC framing -- the two tiers
use different evaluation standards by design (Section~\ref{sec:eval}).}

% =============================================================================
\section{Experimental setup}
% =============================================================================

\subsection{Datasets}
\label{sec:reuse}
Following \cite{daominh2026p0}'s benchmark-role assignment, XES3G5M
\cite{liu2023xes3g5m} (18{,}066 learners, 7{,}652 items, 865 KCs,
6{,}413{,}353 interactions) is the \emph{primary} benchmark: low
consecutive-KC-repeat rate, non-saturated headline AUC, and the clearest
graph-backbone separation among the corpora \cite{daominh2026p0} studied.
ASSISTments~2012 is secondary (single-skill Q-matrix; leakage-ablation
sanity). Junyi Academy provides the ground-truth cross-validation role via
its expert-curated prerequisite annotation.

\subsection{Baselines reused from P0}
We reuse \cite{daominh2026p0}'s published baseline results
(\texttt{baselines/p0\_reused/baseline\_results.csv}, 69 rows, and
\texttt{baseline\_fold\_results.csv}, 205 per-fold rows; both copied
verbatim from \cite{daominh2026p0}'s \texttt{results/tables/}, pinned to
commit \texttt{5e3d6a01c32a649f1655ac8712f03359bddfe478}, vendored
2026-08-02) for BKT, DKT, AKT, simpleKT, GKT, GIKT, SKT, DyGKT, and DGEKT
under train-only, leakage-controlled graph construction, rather than
retraining these nine models on a single RTX~3090. Reuse is valid only if
our preprocessing exactly matches \cite{daominh2026p0}'s (same hashed
\texttt{kc\_id} space, same Q-matrix). \texttt{scripts/01\_verify\_p0\_preprocessing\_match.py}
checks that each DH$^2$A-KT dataset config points at an unmodified,
checksummed copy of \cite{daominh2026p0}'s own config and that the vendored
commit is recorded; for XES3G5M this currently confirms the config file
matches (sha256 \texttt{4f9f35d9\ldots d980b2}) against the vendored commit
above. \todo{This check only verifies the config file is present and
pinned -- it does not yet verify that the locally preprocessed
\texttt{data/} tensors match P0's byte-for-byte (docs/execution-plan.md M0
gate); report that stronger check here once M0's raw data is downloaded and
preprocessed.}

\subsection{Manipulation check}
Following \cite{daominh2026p0}'s Section~4.7 reasoning, we do not interpret
any ablation result as evidence of hypergraph benefit unless DH$^2$-KT first
passes a manipulation check: AUC must move meaningfully under near-total
($p=0.90$) destruction of the constructed hyperedge set. A backbone whose
AUC does not move under total destruction is \emph{graph-inert}, and its
ablation numbers carry no interpretable meaning about graph benefit.
\todo{Report pass/fail and $\Delta$AUC here once
docs/execution-plan.md M6 runs.}

\label{sec:eval}

% =============================================================================
\section{Results}
% =============================================================================
\todo{This entire section is empty by design (docs/execution-plan.md
M5-M9 not yet run at the time this skeleton was drafted). Populate with:
(1) Tier-1 AUC/ACC/NLL vs.\ reused P0 baselines, with 95\% CI; (2) ablation
table (per-node-type, hypergraph-vs-pairwise, static-vs-dynamic,
with/without causal layer); (3) manipulation-check result; (4) ground-truth
cross-validation F1 on Junyi; (5) Tier-2 pilot qualitative table. Do not
fabricate placeholder numbers -- leave this section unwritten until real
results exist, per the anti-pattern warned against in every prior planning
document for this project.}

% =============================================================================
\section{Discussion}
% =============================================================================

\subsection{Threats to validity}
Four risks are known ahead of any experiment and are stated here so they
can be checked against, not discovered after the fact. First, Tier~1's
baseline comparison (Section~\ref{sec:reuse}) is only valid if our
preprocessing exactly reproduces \cite{daominh2026p0}'s hashed
\texttt{kc\_id} space and Q-matrix; \texttt{scripts/01\_verify\_p0\_preprocessing\_match.py}
checks the config level of this today (Section~\ref{sec:reuse}), but the
data-tensor-level check remains open until docs/execution-plan.md
milestone M0's gate is passed. Second, extending hyperedge construction
beyond concept-prerequisite hyperedges may take materially longer than the
session/discussion-thread/teacher-intervention interfaces suggest; if so,
Section~\ref{sec:tier1-graph} already scopes discussion-thread and
teacher-intervention hyperedges to the Tier~2 case study rather than
letting this risk block the Tier~1 quantitative core. Third, any
epoch/batch-budget comparison against \cite{daominh2026p0}'s published
baselines that cannot be matched exactly will be reported as an
\emph{observational} rather than compute-matched comparison, following the
same self-imposed limitation \cite{daominh2026p0} applies to its own
cross-corpus comparisons. Fourth, the manipulation check
(Section~\ref{sec:eval}) is a necessary, not sufficient, condition for
interpreting an ablation result as genuine hypergraph benefit: a backbone
that fails it is graph-inert and its ablation numbers are set aside, but
passing it does not by itself establish that a specific hyperedge kind --
rather than the hypergraph mechanism in general -- is responsible for any
observed gain. \todo{Once M5/M6/M7 results exist, replace this qualitative
framing with the actual outcomes -- did preprocessing verification pass at
the data level, did any risk materialize, what was the compute-budget
match status.}

\subsection{Limitations}
\label{sec:limitations}
\begin{itemize}
  \item Session, discussion-thread, and teacher-intervention hyperedges
  remain data-blocked on all three core benchmarks; Tier~1's quantitative
  claims are scoped to Student-Exercise-Concept(-Hint-Video) unless
  Section~\ref{sec:tier1-graph}'s ES-KT-24 integration is completed.
  \item The causal-effect layer's identification assumption (no unmeasured
  confounding given observed features) is stated, not proven, on
  observational KT data.
  \item The hyperedge leakage audit's pairwise-projection approximation
  (Section~\ref{sec:leakage-ext}) weights all projected pairs equally; a
  hyperedge-native statistic is future work.
\end{itemize}

% =============================================================================
\section{Conclusion and future work}
% =============================================================================
\todo{Write after Results/Discussion are complete.}

\section*{Acknowledgment}
All Tier~1 experiments in this paper are designed to run on a single
consumer-grade NVIDIA RTX~3090 (24GB), matching the compute budget already
validated by \cite{daominh2026p0} on the same benchmarks. \todo{Add funding
source(s) and any human evaluators recruited for the Tier~2 pilot
(Section~\ref{sec:tier2}) once M8/M9 determine whether human evaluation was
collected.}

\bibliographystyle{IEEEtran}
\bibliography{references}

\end{document}
